% sections/experiences.tex
% Expériences Professionnelles — ordre antéchronologique
% Inclut : stages, jobs, engagements associatifs pertinents.
% Chaque bullet utilise la formule PAR ou XYZ.
% Maximum 3 expériences, 2-3 bullets chacune (contrainte 1 page).

\section{Expériences Professionnelles}

% --- Expérience 1 : Oracle (la plus récente et la plus significative) ---
\cventry{06/2025 -- Présent}{Ingénieur IA Stagiaire -- Innovation Lab}{Oracle}{San Luis Potosí, Mexique}{Award : Oracle TrailBlazer Intern 2025}{%
  \begin{itemize}
    \item Architecturé une plateforme UI agentique (LangGraph, A2A, MCP, Oracle ADB, RAG) avec des temps de réponse inférieurs à 30 secondes grâce à un cache sémantique ; présentée lors de plusieurs événements Oracle 2025-2026.
    \item Intégré LangFuse en premier dans le département -- déployé en auto-hébergement sans tierce partie, activement utilisé par 30+ ingénieurs pour l'observabilité des systèmes IA.
    \item Vainqueur du Hackathon interne ECEU/IGIU Memorial 2025 (FieldSync) -- plateforme de synchronisation de données en temps réel pour chantiers de construction.
  \end{itemize}
}

% --- Expérience 2 : LEAD 3D (hardware, fabrication, pont matériel-logiciel) ---
\cventry{01/2025 -- 06/2025}{Stagiaire Ingénieur Conception \& Fabrication}{LEAD 3D}{San Luis Potosí, Mexique}{}{%
  \begin{itemize}
    \item Conçu, modélisé et fabriqué des pièces mécaniques en impression 3D (Bambu Lab, Fusion 360) avec prise en charge complète du flux de production : configuration machine, calibration, maintenance et sélection matériaux.
    \item Réalisé le scan 3D optique (lumière bleue et IR, scanner Raptor) de composants industriels réels (chaise 1:1, machine à café, jante de golf) pour ingénierie inverse ; post-traitement des nuages de points sous Blender.
  \end{itemize}
}

% --- Expérience 3 : Nexon Automation / LamBot FRC (robotique, leadership, résultats mondiaux) ---
\cventry{06/2022 -- 12/2023}{Mentor STEM \& Capitaine Division Ingénierie}{Nexon Automation -- Équipe LamBot 3478 FRC (alliée General Motors)}{San Luis Potosí, Mexique}{}{%
  \begin{itemize}
    \item Dirigé le développement logiciel de 2 robots compétitifs de niveau mondial en 10+ événements : Semifinaliste Division Mondiale FRC 2023, Top 5 Division Mondiale FRC 2022.
    \item Développé 2 modèles de détection d'objets en temps réel (OpenCV, réseaux de neurones, Haar cascade) pour contrôle autonome, et appliqué des contrôleurs PID et de profilage de mouvement pour la cinématique du bras (saison 2023).
    \item Enseigné la POO (Java, Python, WPILib) à 2 groupes de 60 étudiants ; encadré les domaines électrique et programmation de l'équipe.
  \end{itemize}
}

% TODO: Ajouter Scale AI (01/2024 -- 08/2024) si l'espace le permet après compilation.
%       Entrée compacte : "Entraîneur de modèles IA, Scale AI — Évaluation et débogage de 5 modèles de génération de code (Python, Java, C++, Shell) ; développement de datasets d'entraînement pour 2 modèles spécialisés."
% TODO: Ajouter lettre de mission LEAD 3D [alejandro] au dossier Campus France.
